\documentclass[a4paper,12pt]{article}
\usepackage[left=1cm,right=1cm,top=1cm,bottom=1cm]{geometry}
\usepackage{chemfig}
\usepackage[version=4]{mhchem}
\usepackage{mol2chemfig}
\usepackage{tabularx}

\title{Correction du TP2}
\author{Yannick Patois}
\date{12/11/2025}

\begin{document}
\maketitle

\section*{ACTIVITÉ 1 : des représentations pour comprendre la notion de chiralité}

\subsection*{1.Pour chacune des molécules suivantes, donner la représentation de Cram et indiquer si la molécule est chirale.}

La première molécule représentée est un 1-chloropropan-1,2-ol \ce{CHClOH-CH2OH-CH3}.
Une représentation de Cram de cette molécule est~:

\vfill
\begin{center}
    \chemfig[bond join=true]{
        {\color{red}*}(-[:0]{\color{red}*}(-[:-30]CH3)(<:[:60]H)(<[:30]OH))%
        (-[:150]Cl)(<[:-120]OH)(<:[:-150]H)
    }\\
    1-chloropropan-1,2-ol \ce{CHClOH-CH2OH-CH3}.
\end{center}
\vfill

La seconce molécule représentée est un 1,2-dichloropropan-1,2-ol \ce{CHClOH-CH2ClOH-CH3}.
Une représentation de Cram de cette molécule est~:

\vfill
\begin{center}
    \chemfig[bond join=true]{
        {\color{red}*}(-[:0]{\color{red}*}(-[:-30]CH3)(<:[:60]Cl)(<[:30]OH))%
        (-[:150]Cl)(<[:-120]OH)(<:[:-150]H)
    }\\
    1,2-dichloropropan-1,2-ol \ce{CHClOH-CH2ClOH-CH3}.
\end{center}
\vfill

\subsection*{Pour chacune des molécules suivantes, construire la molécule à l’aide des modèles moléculaires et compléter le
tableau ci-dessous.}

\begin{table}[h]
    \centering
    \begin{tabularx}{\textwidth}{|X|c|c|c|X|}
    \hline
    \textbf{Molécule} & \textbf{Groupe} & \textbf{Famille} & \textbf{Chirale ?} & \textbf{Justification} \\
    \hline
    3-chlorobutan-2-ol                 & Alcool    & Alcool             & Oui & Deux carbones asymétriques \\
    \hline
    Acide 3-hydroxy-2-methylbutanoique & Carboxyle & Acide carboxylique & Oui & Deux carbones asymétriques \\
    \hline
    2-methylbutanal                    & Carbonyle & Aldhéyde           & Oui & Un carbone asymétrique \\
    \hline
    Ethanoate de propyle               & (Ester)   & (Ester)            & Non & Deux hydrogènes \\
    \hline
    \end{tabularx}
\end{table}

\end{document}
